\documentclass[11pt]{article}
\usepackage{amsmath}
\title{A qualified release and a safe timer can jointly shorten an authority handoff}
\date{26 November 2026}
\begin{document}
\maketitle
\paragraph{Law and evidence.}
A shared actuator is subject to a modeled exclusive-effect law:
the successor must not start while residual effects of the old
operator can still occur. The old operator withdraws at local
time zero. Its first effect-free cut is $-\delta_A+q$ in
common time for old clock offset $\delta_A$ and tail $q\ge0$.
In the declared horizon the local clocks have equal rates and
constant offsets. The new operator has offset $\delta_B$ and
a separately granted next authorization, effective when either
of the permitted transitions below starts. It gets no live
advance observation of the old operator except a possible
qualified release. Such a release is bound to the same actor,
effect resource and command version and cannot arrive before
the old effect-free cut. Its latency after that cut is $\lambda$,
or it is absent. Scoped delay intervals describe admitted
joint traces, not measured distributions or an implementation
of trusted release. Neither a task pause, stop request receipt,
nor a software acknowledgement that retirement started qualifies
by itself. Establishing an actual effect-free cut is external.

\paragraph{Exact fixed-fallback comparison.}
A timer at successor local time $x$ would start in common time
$x-\delta_B$. Put $d=\delta_B-\delta_A+q$.
When trying the first of the safe timer and a qualified release,
the effect-free blackout is $\min(x-d,\lambda)$ on a trace
with release and $x-d$ otherwise. A negative value represents
overlap. Because $\lambda\ge0$, no-overlap on every trace
is equivalent to $x\ge\max d$. A finite nonempty set of
joint traces is assumed. A later timer cannot shorten any
trace's blackout, so $x=\max d$ is optimal among these
fixed-fallback policies, with or without qualified release.
This is not an optimality claim for policies with extra sensors,
negotiation, retry or varying authority. The Lean kernels check
the quantified no-overlap equivalence and monotonicity of the
pointwise blackout. A direct world-time enumeration confirms
3{,}200 cases across clock, tail and receipt variations.

\paragraph{Two witnesses.}
Consider traces $(d,\lambda)=(1,1)$ and $(5,4)$.
The safe timer alone leaves blackouts 4 and 0; a release-only
policy leaves 1 and 4. Combining them gives 1 and 0, hence
a one-tick worst case instead of four for each isolated policy.
Both sources of evidence must be available under the same declared
authorization. Pairing the old early cut with the slow receipt
from the other hypothesis creates a four-tick case, as would
dropping the early receipt in the first actual trace. Reusing a
timer certified with zero tail where the old tail is now one
can produce a one-tick overlap even though the timestamps still
look valid. Equality is allowed precisely at first effect-free
cut. A noncausal or stale receipt could change that safety
conclusion and is outside the model.

\paragraph{Scope.}
The two decision channels are combined on the same joint trace,
not by min of incompatible worst cases. Logical grants can
overlap or be absent; this model's exclusive actuation law
is about effective control of one shared effect, not the number
of task holds or a general owner election. No controller real
time, physical validity of a receipt, or closed architecture item
is claimed. The handoff gives a scoped F26 comparison and a
producer requirement for a qualified release interface.
\end{document}
